\section{Topic, Scope and Supervision}

\subsection{Chosen Topic}

The chosen topic is \textbf{memory consistency models} in shared-memory
multiprocessors.

\medskip
\noindent\textbf{Working title.}\\
\emph{Revisiting Sequential Consistency: Memory Consistency Models, Their Cost,
and Whether Speculation Has Made Strong Ordering Affordable}

\medskip
\noindent\textbf{Scope.}
A shared-memory multiprocessor must specify when, and in what order, one core's
writes become visible to the others.
That specification is the machine's memory consistency model, and it is an
architectural contract as binding as the instruction set.
The paper covers the design space from sequential consistency through TSO, PSO,
weak ordering and release consistency; the mechanisms that make relaxation
necessary (store buffers, out-of-order execution) and those that could make it
unnecessary (speculative execution with coherence-based detection); and the
historical argument that led the industry to reject sequential consistency
around 1990.

\medskip
\noindent\textbf{Guiding question.}
Sequential consistency was rejected as a hardware memory model on performance
grounds in the late 1980s, on machines that did not speculate.
A modern out-of-order core already contains both mechanisms a strong model
requires --- coherence invalidations that announce remote writes, and a
squash-and-restart path built for branch misprediction recovery.
The paper asks whether the performance objection still holds, and, if it does
not, what else accounts for the continued absence of sequentially consistent
hardware.

\medskip
\noindent The work is organised around four questions:

\begin{description}[leftmargin=2.2cm,style=nextline,font=\normalfont\bfseries]
\item[Q1 Formal] Can SC, TSO, PSO, weak ordering and release consistency be
  expressed as successively weaker constraints on a single set of ordering
  relations, with strict containment proved rather than asserted?
\item[Q2 Empirical] Which relaxations are \emph{observable} on real x86 and ARM
  hardware, as distinct from those merely \emph{permitted} by the architecture?
\item[Q3 Quantitative] Holding the microarchitecture fixed and varying only the
  memory model, what does sequential consistency cost with and without
  speculative load execution, as core count rises?
\item[Q4 Interpretive] If speculative SC proves inexpensive, what accounts for
  its absence from shipping processors?
\end{description}

\subsection{Mapping to the Prescribed Textbook}

The prescribed text is S.\,R.~Sarangi, \emph{Basic Computer Architecture},
version 3.09 \cite{sarangi2025basic}.
The topic maps principally to \textbf{Chapter 12, \emph{Multiprocessor
Systems}}, and within it to \textbf{\S12.4, \emph{MIMD Multiprocessors}}.

\begin{table}[H]
\centering
\small
\begin{tabular}{@{}llp{6.6cm}@{}}
\toprule
\textbf{\S} & \textbf{Title (page)} & \textbf{Role in the paper} \\
\midrule
\multicolumn{3}{@{}l}{\textbf{Primary --- Chapter 12}} \\
12.4.1 & Logical Point of View (579)          & Shared-memory abstraction. \\
12.4.2 & Coherence (581)                      & The coherence/consistency distinction, developed at length. \\
12.4.3 & \textbf{Memory Consistency (583)}    & \textbf{The core section. The paper is an extension of this.} \\
12.4.5 & Shared Caches (596)                  & Background. \\
12.4.6 & Coherent Private Caches (596)        & Extended to full MESI/MOESI and directory protocols. \\
12.4.7 & Implementing a Memory Consistency Model$^*$ (603) & Direct antecedent of the speculative-SC discussion. \\
\addlinespace
\multicolumn{3}{@{}l}{\textbf{Supporting}} \\
10.4   & Pipeline Hazards (420)               & Reordering and dependences. \\
10.9   & Performance Metrics (463)            & Benchmarking framework. \\
10.11.4& Out-of-Order Pipelines (492)         & Speculation and the squash mechanism. \\
11.2   & Caches (517)                         & Private caches, non-blocking misses. \\
11.3   & The Memory System (532)              & Latency figures motivating the store buffer. \\
12.2.2 & Shared Memory vs Message Passing (571)& Framing. \\
12.2.3 & Amdahl's Law (576)                   & Scaling discussion. \\
12.6   & Interconnection Networks (620)       & Coherence latency in the simulator. \\
App.~A & Cortex-A8/A15, Sandy Bridge (734--746)& Real-machine case studies. \\
\bottomrule
\end{tabular}
\caption{Mapping of the topic onto the prescribed text. $^*$ denotes a section
marked advanced.}
\end{table}

\noindent\textbf{Additional texts.} Because \S12.4.3 treats only sequential
consistency and the relaxed models are outside the scope of the prescribed
text, three further sources carry the technical weight of the paper.
All are fully specified, with reproduced tables of contents, in
Chapter~7 of the paper draft.

\begin{itemize}[leftmargin=*,itemsep=2pt]
\item S.\,R.~Sarangi, \emph{Next-Gen Computer Architecture: Till the End of
  Silicon}, White Falcon 2023, ISBN 8119510143 \cite{sarangi2023nextgen}.
  Chapter~9, \emph{Multicore Systems: Coherence, Consistency and Transactional
  Memory}, runs to roughly 118 pages and is the source of the formalism used
  throughout.
\item V.~Nagarajan, D.\,J.~Sorin, M.\,D.~Hill and D.\,A.~Wood, \emph{A Primer
  on Memory Consistency and Cache Coherence}, 2nd ed., Springer 2020
  \cite{nagarajan2020primer}. The standard reference on the subject.
\item Vendor architecture manuals (Intel SDM Vol.~3A; Arm ARM Ch.~B2;
  RISC-V RVWMO), cited in preference to any textbook wherever the paper states
  what a particular instruction set guarantees.
\end{itemize}

\subsection{Advisor}

I will consult \textbf{Prof.~[Sumantra Dutta Roy / Subrat Kar --- SELECT ONE]}
as advised.

\medskip
\noindent Zotero library \texttt{ELL305-MemoryConsistency} has been created and
will be shared with Prof.~Subrat Kar. Bibliography is exported to
\texttt{refs.bib} and rendered via BibTeX.

